\section{A neural implementation to test}
\label{sec:architecture}

\subsection{Selected measurements before the first destructive operation}

At boundary \(\ell\), define \(h_\ell\) to be the state \emph{before} the
operation that could erase the chosen distinction. The ordinary path is
\[
 h_{\ell+1}=F_\ell(h_\ell),
\]
for example \(F_\ell(h_\ell)=\rho(P_\ell h_\ell)\). A candidate protected
record is
\begin{equation}
 r_\ell=\CReLU(U_\ell h_\ell).
 \label{eq:selectivecrelu}
\end{equation}
The selected signed measurement is recovered by
\[
 U_\ell h_\ell=(r_\ell)_+-(r_\ell)_-.
\]
A projection with \(k\) independent rows creates \(k\) signed measurements
and \(2k\) nonnegative CReLU channels. Count both correctly when comparing
activation budgets.

Because the record branches before \(F_\ell\), it can retain a contrast that
the ordinary projection would remove. If \(h_\ell=Px\), every subsequent
\(U_\ell h_\ell\) is identical on inputs having the same \(Px\). It cannot
separate an admissible pair that differs in a target-relevant nullspace
direction unless another retained observation already distinguishes that pair.
Record placement must precede the first
relevant loss, rather than merely preceding the final rectifier.

The retained collection can be written
\[
 E(x)=\left(h_L(x),
 \{(\ell,k_\ell,r_\ell(x)):\ell\in\mathcal B\}\right),
\]
where \(k_\ell\) is a chart or branch identity when needed and
\(\mathcal B\) is the set of protected boundaries. Known shared transition
parameters need not be copied into every input record. Input-dependent
selection and map choices must be interpretable by the decoder.

This protects selected measurements, not automatically the target. The
requirement at a boundary remains
\[
 (F_\ell(h),r_\ell(h))=(F_\ell(h'),r_\ell(h'))
 \Longrightarrow T_\ell(h)=T_\ell(h'),
\]
including other persistent records whenever \(T_\ell\) depends on them.
An exact correction quantity can sometimes replace selected signed
measurements, as in the ReLU example. A learned approximate correction needs
its own error analysis.

\subsection{An end-to-end linear audit}

For a linear chain \(h_{\ell+1}=A_\ell h_\ell\), \(h_0=x\), with records
\(r_\ell=N_\ell h_\ell\), the final observation matrix is
\begin{equation}
 \mathcal M_L=
 \begin{bmatrix}
 A_{L-1}\cdots A_0\\
 N_0\\
 N_1A_0\\
 \vdots\\
 N_{L-1}A_{L-2}\cdots A_0
 \end{bmatrix}.
 \label{eq:chain}
\end{equation}
For original target \(C_0x\), the complete criterion is
\[
 \ker\mathcal M_L\subseteq\ker C_0.
\]
This catches loss at later aggregation even when a record was adequate when
first emitted. Nonlinear versions use the global fiber criterion from
Theorem~\ref{thm:fibers}.

\subsection{Layer order and heavier upper weights}

Layer indices, transition direction, and map domains are part of the object
being glued. Composition is ordered. Heavier upper weights can be used as
known positive scaling of individual records or as reconstruction weights.
Nonzero scaling can be undone while the individual records remain available.

A single weighted scalar total is another compression and requires its own
proof. Distinct weights alone are insufficient:
\[
 1(1)+2(3)+3(2)=13=1(2)+2(1)+3(3).
\]
Even the same values in different orders collide in this example.

The role of weights must be declared. Position labeling, confidence in noisy
measurements, and semantic weighting of the target are different operations.
The experiments must not alter an unweighted target while claiming merely
to improve its encoding. Upper-layer weights can express the intended
priority, while the retained indexed records supply the information needed
for exact interpretation.

\subsection{Learning proposal}

A candidate objective is
\begin{align}
 \mathcal L={}&\mathcal L_{\rm task}
 +\lambda_{\rm con}\sum_{e=(i,j)}
   \norm{\pi_{i,e}(v_i)-\pi_{j,e}(v_j)}^2 \nonumber\\
 &+\lambda_{\rm cyc}\sum_{(i,j,k)}
   \norm{\tau_{ik}(v_i)-\tau_{jk}(\tau_{ij}(v_i))}^2
 +\lambda_{\rm budget}\mathcal B(U).
 \label{eq:training}
\end{align}
The cycle term applies only on valid common domains when these transition
maps are part of the model. The budget term can penalize retained dimensions
or grouped projection rows. This is a proposed training procedure, not a
preservation theorem. Learned zero maps can make consistency losses vanish;
hard sparsification can delete a necessary record; and a finite penalty need
not impose an exact constraint.

Start with a fixed record budget and no input-dependent hard deletion.
After the basic mechanism works, test any adaptive gate for target-changing
collisions in regions where protection turns off. Oracle teacher directions
may be used in explicitly labeled diagnostic controls. Ordinary training
should use the declared data and task loss rather than hidden teacher
features.

Inference must remain causal: records use current and earlier states.
Future-layer supervision can be used in training under a specified protocol,
but evaluation must not reveal unavailable future observations. This
implements the causal discipline emphasized in the research discussion
\cite{discussion}.

\subsection{Forward values, gradients, and arithmetic}

An exact forward identity implies equal ordinary derivatives wherever the
functions are differentiable, provided the computation remains connected.
Detaching or approximating a correction can change learning even when some
forward values match. At internal activation boundaries, equivalent expressions
can produce different automatic derivatives, which need not equal the
derivative of the complete function.

For example, \(\rho(z)-\rho(-z)=z\) as a real function. If both ReLU branches
use derivative zero at \(z=0\), naive differentiation of that graph returns
zero there, whereas the identity function has derivative one. A matched
backward convention is needed when gradient preservation is part of a claim.

For the pair correction, the absolute-value subtraction can also suffer
catastrophic cancellation. At \(A=-10^{20},B=1\), its mathematical value is
one, but the subtraction formula can return zero in binary64. A more stable
pair evaluation returns zero when signs match and
\(\min(|A|,|B|)\) when signs oppose. Compare signs directly rather than
forming \(AB\), which may overflow. This improves that local computation;
it does not make the complete floating-point pipeline exact.

A local diagnostic
\[
 \ker DE(s)\subseteq\ker DT(s)
\]
is necessary for a differentiable decoder factorization but is not a
global proof. Distinct activation regions may collide. The map \(x\mapsto
x^2\) is locally invertible on separate positive and negative intervals,
yet recovering \(x\) requires a branch identity. A neural atlas must
therefore address cross-region collisions and valid chart domains.

\section{Evaluation plan and acceptance criteria}
\label{sec:experiments}

\subsection{Stage 0: mathematical contracts}

Reproduce the identities and negative controls using the supplied deterministic
verification script. Build a small linear prototype with known observation and
target matrices, construct a minimum residual basis, and verify its explicit
decoder with rational arithmetic. This checks the mathematical implementation;
it is not a neural-learning result.

The negative controls are essential. Merged CReLU must fail on a missing
contrast. Clipped-sum-plus-contrast must fail on its stated pair. The stated
pairwise-product record must fail to determine the triple-product target. Floating-point tests
must exhibit the stated rounding and subtraction failures. Favorable
examples alone would not distinguish the proposed contract from a weaker one.

\subsection{Stage 1: the original small-network diagnosis}

Reuse the original teacher family and, where accessible, saved failed
checkpoints. First freeze the ordinary representation and train only the
selective measurement and readout components. Then run an end-to-end
comparison separately. Keep training, model-selection, and final evaluation
data distinct, with matched teachers and initializations across variants.

Separate failure at the linear projection from failure at ReLU. An augmented
span floor can certify remaining linear information loss. A zero span floor
does not certify a hybrid clipped representation; activation-collision
checks must follow. The QDO derivation appendix establishes the relevant
exact floor and lower-bound distinction.

Use the teacher-aware correction only as an oracle positive control. It
demonstrates that a suitable record can repair a chosen failure. It does
not establish that a learned model discovers that record from examples.

\subsection{Comparators and resource accounting}

\begin{longtable}{@{}>{\raggedright\arraybackslash}p{0.30\linewidth}>{\raggedright\arraybackslash}p{0.64\linewidth}@{}}
\toprule
Comparator & Question it answers\\
\midrule
\endhead
Ordinary ReLU model & What are the unmodified target error and
collision patterns?\\
CReLU on the same projected responses & How much loss comes from
sign clipping after the existing measurement?\\
Ordinary wider model & Is the effect explained by added capacity or compute?\\
Random protected measurements & Does learned selection outperform a matched
random record budget?\\
Unstructured skip measurements & Does atlas alignment add value beyond
carrying the same amount of extra data?\\
Selected records without learned transport & Are the transition maps
responsible for any gain?\\
Uniform versus increasing layer weights & Does upper-layer weighting help
after identity and capacity are held fixed?\\
Full-rank signed measurement reference & What does a broad,
input-preserving reference achieve? Report its larger budget explicitly.\\
Oracle contrast or correction record & Can the diagnosed failure be repaired
when the missing quantity is supplied? Keep this separate from learned results.\\
\bottomrule
\end{longtable}

Match parameter count, activation count, persistent-state bytes, and measured
compute as closely as possible. These budgets may not all match simultaneously;
report the differences and use separate capacity-matched and compute-matched
comparisons where needed. A low-dimensional dense projection can still be
computationally expensive.

\subsection{Stage 2: multiple boundaries and reusable targets}

Test deeper compositions with records emitted before multiple lossy boundaries.
Include order-sensitive targets, layer permutations, chart relabeling with
correctly transformed maps, delayed dependencies, and mixed-sign coefficients.
Include higher-order targets that cannot be determined by the selected
pairwise records without an anchor.

Use a declared collection of readouts on the same underlying features to test
reuse. A record that stores one scalar answer can perform perfectly on that
answer while failing a broader requirement. Distinguish interpolation on
the training distribution from the stated admissible domain.

\subsection{Stage 3: conditioning and numerical behavior}

Perturb records along target-visible singular directions and compare target
error with the derived stability bound. Include consistent perturbations
\(\varepsilon=Mv\): they expose target errors a compatibility residual
cannot detect.

Report precision, scaling, overflow handling, and whether correctness means
functional tolerance or bitwise equality. Test cancellation near zero and
widely separated magnitudes. Check gradient flow through the protected path
and state the subgradient convention at activation boundaries.

\subsection{Metrics and decision gates}

\begin{enumerate}
\item \textbf{Structural correctness.} For analytic prototypes, prove the
fiber or kernel condition on the declared domain. For learned nonlinear
models, report explicit witnesses or bounded-domain verification where
tractable. Failure to find a collision is not proof of sufficiency.
\item \textbf{Target performance.} Report held-out normalized squared error
for the original task and appropriate metrics for added tasks. The original
\(10^{-3}\) NMSE threshold can be retained for continuity, but it is not an
exact-preservation criterion.
\item \textbf{Mechanism.} Audit the identified contrasts, negative common
components, chart consistency, branch identities, anchors, and survival of
records through the complete computation.
\item \textbf{Cost.} Report parameters, FLOPs, latency, bytes, and record
count alongside accuracy. Attribute benefits to selectivity or gluing only
after the appropriate controls.
\item \textbf{Reproducibility.} Preserve seeds, data splits, checkpoint
selection rules, arithmetic settings, and per-case results. Report
uncertainty across independent runs for learned comparisons.
\end{enumerate}

Proceed to a larger model only if learned records repair the identified
failures at a defensible budget and the gain survives ordinary-capacity,
random-record, and unstructured-skip controls. If agreement improves while
target ambiguity persists, reject that representation. If a simple added
channel explains the gain, report it without attributing the gain to atlas
structure.

\section{Established results and open questions}

Selective preservation is mathematically possible: a representation may lose
source distinctions while exactly retaining a declared target. The ReLU
correction gives a constructive nonlinear example. The finite-pile identity
shows that local corrections can compose. The linear theorems identify the
missing directions requiring side records and give an explicit decoder.
The atlas formulation separates compatibility, source explanation, and target
observability.

The unresolved research question is whether a neural system can learn useful
records and transport maps at a better accuracy--cost tradeoff than simpler
baselines. Selection must occur before collapse, preserve required nonlinear
branches, and remain sufficient for later computation. Upper-layer weights,
successful prediction, small consistency loss, and full parameter-Jacobian
rank do not remove those obligations.

The first implementation should be a small forward-boundary module with a
declared record budget, a reconstruction interface to the atlas, and evaluation
against the original certified failure types. Its acceptance statement is:
\begin{keybox}
Every ambiguity left by the complete retained state is invisible to the
declared task family, and the implemented reconstruction exposes that
information to the readout at the stated numerical precision.
\end{keybox}

\clearpage
\section{Reproducible mathematical checks}
\label{sec:checks}

The source bundle includes \code{verify_math.py} and its generated
\code{verification_results.json}. The checks use exact integer and rational arithmetic
for algebraic identities and explicit floating-point operations only for
rounding counterexamples. All eight check groups passed.

\begin{center}
\small
\begin{tabularx}{\linewidth}{@{}p{0.47\linewidth}X@{}}
\toprule
Check & Domain and result\\
\midrule
Two-input correction & 289 pairs in \([-8,8]^2\); absolute, contrast,
gate, and reconstruction identities agree.\\
Signed coefficients & 2,401 coefficient/input quadruples in
\([-3,3]^4\); general correction agrees with the target.\\
Positive position weights & 294 ordered weighted cases; positive-homogeneity
and weighted correction checks pass.\\
Original four-feature teacher & 2,401 feature tuples in \([-3,3]^4\);
the pair corrections recover the target exactly.\\
Finite piles & 3,125 five-leaf stacks in \([-2,2]^5\), each through
all 14 ordered binary groupings: 43,750 tree evaluations.\\
Deliberate ambiguities & CReLU after a merge, clipped sum plus contrast,
and pairwise-product parity retain their target-changing collisions.\\
Selective linear record & 625 inputs in \([-2,2]^4\); one added coordinate
repairs the required target while an irrelevant coordinate stays unrecoverable.\\
Finite precision & float32/64 rounding collisions and subtractive
cancellation are reproduced; the stable pair formula is checked.\\
\bottomrule
\end{tabularx}
\end{center}

These finite checks validate calculations and witnesses. The proofs establish
the general statements under their assumptions. No new network training,
measured speedup, or accuracy improvement is reported.
